\documentclass[aspectratio=169,10pt]{beamer}
\usetheme{metropolis}
\usepackage{booktabs,amsmath,amssymb}
\title{Q07: The Science of Scaling}
\subtitle{Power laws, Chinchilla, and the limits}
\author{DLS 2026 --- Kukla}
\date{October 2026}
\begin{document}
\maketitle

\begin{frame}{The question}
\begin{block}{Student question (\#21, \#22, \#39)}
Why does scaling work? Is there a ``science of scaling,'' what are its limits, and can learning theory
model it?
\end{block}
\end{frame}
\begin{frame}{The empirical finding}
\[ L(N) \approx L_\infty + (N_0/N)^{\alpha} \]
\begin{itemize}
\item Test loss falls as a \textbf{power law} in size $N$, data $D$, compute $C$.
\item Holds over many orders of magnitude.
\item Lets labs \emph{predict} loss \emph{before} training.
\end{itemize}
\end{frame}
\begin{frame}{Compute-optimal (Chinchilla)}
For fixed $C \approx 6ND$: scale $N$ and $D$ \textbf{together}.
\begin{itemize}
\item Earlier models: over-parameterized, under-trained.
\item Smaller model + more data can beat bigger model + less, at equal cost.
\end{itemize}
\alert{Scaling laws are prescriptive, not just descriptive.}
\end{frame}
\begin{frame}{Why would it work?}
Partial accounts (no closed theory):
\begin{itemize}
\item Prediction as \textbf{compression} (lower loss = better world-model).
\item \textbf{Data-manifold} dimension sets the exponent.
\item \textbf{Over-parameterization} + SGD find generalizing solutions.
\end{itemize}
\end{frame}
\begin{frame}{The limits}
\begin{itemize}
\item \textbf{Data wall}: high-quality human text is finite.
\item \textbf{Diminishing returns}: small exponent $\Rightarrow$ exponential cost.
\item \textbf{Loss $\ne$ capability}; emergence debated.
\item \textbf{New axes}: post-training + inference-time compute.
\end{itemize}
\end{frame}
\begin{frame}{Can theory explain it?}
\begin{itemize}
\item Classical bounds (VC) predict overfitting --- too loose.
\item Data-manifold / solvable models recover power laws under assumptions.
\item No predictive theory matches real exponents across architectures yet.
\end{itemize}
\end{frame}

\end{document}
